\documentclass[11pt,letterpaper]{article}
\usepackage[margin=0.85in]{geometry}
\usepackage{amsmath,amssymb}
\usepackage[T1]{fontenc}
\usepackage{lmodern}
\usepackage{fancyhdr}
\usepackage{hyperref}
\pagestyle{fancy}\fancyhf{}
\fancyhead[L]{ENGR 202\quad Lab 5}\fancyhead[R]{Group 35}
\fancyfoot[C]{\thepage}
\setlength{\headheight}{14pt}
\setlength{\parindent}{0pt}
\setlength{\parskip}{6pt}
\newcommand{\Vs}{\widetilde V_s}
\newcommand{\Vz}{\widetilde V_z}
\newcommand{\Vk}{\widetilde V_{1k}}
\begin{document}
\begin{center}
{\Large\bfseries Lab 5 Analytical Predictions}\\[5pt]
Mystery Box Challenge Part 1\\
Tong Su, Carter Flynn Swanson, Zachary Kramer
\end{center}
\section*{Circuit and phasor conventions}
The supplied schematic has a known resistor $R_0=1000\,\Omega$ in series
with the mystery-box impedance $Z$. Channel 2 measures the source voltage;
Channel 1 measures the voltage across the mystery box. Thus
\[
\Vs=\widetilde V_{\mathrm{Ch2}},\qquad
\Vz=\widetilde V_{\mathrm{Ch1}},\qquad
\Vk=\widetilde V_{\mathrm{Ch2}}-\widetilde V_{\mathrm{Ch1}}.
\]
Use the cosine convention $v(t)=\operatorname{Re}\{\widetilde V e^{j\omega t}\}$,
with $\omega=2\pi f>0$. Phasor magnitudes are \textbf{peak amplitudes}, not RMS:
$\Vs=1\angle0^\circ\,\mathrm V$, corresponding to $2\,\mathrm{V_{pp}}$.
The components are ideal and their values are positive. Ohm's law gives
\begin{align}
\widetilde I&=\frac{\Vs}{R_0+Z}, &
\Vz&=\Vs\frac{Z}{R_0+Z}, &
\Vk&=\Vs\frac{R_0}{R_0+Z}. \label{eq:divider}
\end{align}
\section*{Q3.1.1 Unknown resistor}
For a mystery-box resistor $R$, its impedance is $Z=R$. Therefore
\[
\widetilde I=\frac{1\,\mathrm V}{R_0+R},\qquad
\Vz=\frac{R}{R_0+R}\,\mathrm V,\qquad
\Vk=1\,\mathrm V-\frac{R}{R_0+R}\,\mathrm V
       =\frac{R_0}{R_0+R}\,\mathrm V.
\]
All ratios are positive real numbers, so each voltage has zero phase:
\begin{align*}
\boxed{\Vs=1\angle0^\circ\,\mathrm V}\qquad &\text{Channel 2},\\[7pt]
\boxed{\Vz=\widetilde V_R=\frac{R}{R_0+R}\angle0^\circ\,\mathrm V}
&\qquad\text{Channel 1},\\[7pt]
\boxed{\Vk=\frac{R_0}{R_0+R}\angle0^\circ\,\mathrm V}
&\qquad\text{Channel 2 minus Channel 1}.
\end{align*}
Here $R_0=1000\,\Omega$. Both voltages are in phase with the source,
and their sum is exactly $1\angle0^\circ\,\mathrm V$.

\newpage
\section*{Q3.1.2 Unknown capacitor}
For a mystery-box capacitor $C$,
\[
Z_C=\frac{1}{j\omega C},\qquad x=\omega R_0 C>0.
\]
Substituting into the voltage divider and multiplying numerator and
denominator by $j\omega C$ gives
\begin{align*}
\Vz=\widetilde V_C
&=(1\,\mathrm V)\frac{1/(j\omega C)}{R_0+1/(j\omega C)}
=\frac{1}{1+jx}\,\mathrm V\\
&=\frac{1-jx}{1+x^2}\,\mathrm V.
\end{align*}
Using the complex conjugate,
\[
|\Vz|^2=\Vz\Vz^*
=\frac{(1-jx)(1+jx)}{(1+x^2)^2}\,\mathrm V^2
=\frac{1}{1+x^2}\,\mathrm V^2.
\]
The real part is positive and the imaginary part is negative. Hence
\[
|\Vz|=\frac{1}{\sqrt{1+x^2}}\,\mathrm V,\qquad
\arg\Vz=-\tan^{-1}x.
\]
The required phasors are
\[
\boxed{\Vs=1\angle0^\circ\,\mathrm V},\qquad
\boxed{\Vz=\widetilde V_C=
\frac{1}{\sqrt{1+(\omega R_0 C)^2}}
\angle\!\left[-\tan^{-1}(\omega R_0 C)\right]\,\mathrm V}.
\]
\subsection*{Extra credit Voltage across the known resistor}
Subtract the \emph{complex phasors}, not their magnitudes:
\begin{align*}
\Vk&=\Vs-\Vz
=\left(1-\frac{1}{1+jx}\right)\mathrm V
=\frac{jx}{1+jx}\,\mathrm V
=\frac{x^2+jx}{1+x^2}\,\mathrm V,\\
|\Vk|^2&=\frac{x^4+x^2}{(1+x^2)^2}\,\mathrm V^2
=\frac{x^2}{1+x^2}\,\mathrm V^2,\\
\arg\Vk&=90^\circ-\tan^{-1}x.
\end{align*}
Therefore, with $R_0=1000\,\Omega$,
\[
\boxed{\Vk=\widetilde V_R=
\frac{\omega R_0 C}{\sqrt{1+(\omega R_0 C)^2}}
\angle\!\left[90^\circ-\tan^{-1}(\omega R_0 C)\right]\,\mathrm V}.
\]
The capacitor voltage lags the source; the known-resistor voltage leads it.
Their phases differ by $90^\circ$. In these polar answers, inverse tangents
are evaluated in degrees. At low frequency, $\Vz\to\Vs$ and $\Vk\to0$;
at high frequency, $\Vz\to0$ and $\Vk\to\Vs$.

\newpage
\section*{Q3.1.3 Unknown inductor}
For a mystery-box inductor $L$,
\[
Z_L=j\omega L,\qquad y=\frac{\omega L}{R_0}>0.
\]
The voltage across the inductor is
\begin{align*}
\Vz=\widetilde V_L
&=(1\,\mathrm V)\frac{j\omega L}{R_0+j\omega L}
=\frac{jy}{1+jy}\,\mathrm V\\
&=\frac{jy(1-jy)}{1+y^2}\,\mathrm V
=\frac{y^2+jy}{1+y^2}\,\mathrm V.
\end{align*}
Its magnitude and phase follow from
\begin{align*}
|\Vz|^2&=\frac{(y^2+jy)(y^2-jy)}{(1+y^2)^2}\,\mathrm V^2
=\frac{y^2}{1+y^2}\,\mathrm V^2,\\
\arg\Vz&=\arg(jy)-\arg(1+jy)
=90^\circ-\tan^{-1}y.
\end{align*}
Thus
\[
\boxed{\Vs=1\angle0^\circ\,\mathrm V},\qquad
\boxed{\Vz=\widetilde V_L=
\frac{\omega L}{\sqrt{R_0^2+(\omega L)^2}}
\angle\!\left[90^\circ-\tan^{-1}\!\left(\frac{\omega L}{R_0}\right)\right]\,\mathrm V}.
\]
The equivalent positive phase is $\tan^{-1}(R_0/(\omega L))$.
The worksheet's $V_C$ label in this question is interpreted as $V_L$,
since the assumed component is an inductor.
\subsection*{Extra credit Voltage across the known resistor}
\begin{align*}
\Vk&=\Vs-\Vz
=\left(1-\frac{jy}{1+jy}\right)\mathrm V
=\frac{1}{1+jy}\,\mathrm V
=\frac{1-jy}{1+y^2}\,\mathrm V,\\
|\Vk|&=\frac{1}{\sqrt{1+y^2}}\,\mathrm V,\qquad
\arg\Vk=-\tan^{-1}y.
\end{align*}
With $R_0=1000\,\Omega$,
\[
\boxed{\Vk=\widetilde V_R=
\frac{R_0}{\sqrt{R_0^2+(\omega L)^2}}
\angle\!\left[-\tan^{-1}\!\left(\frac{\omega L}{R_0}\right)\right]\,\mathrm V}.
\]
The inductor voltage leads the source; the known-resistor voltage lags it.
Inverse tangents in the polar answers are in degrees. At low frequency,
$\Vz\to0$ and $\Vk\to\Vs$; at high frequency,
$\Vz\to\Vs$ and $\Vk\to0$.

\textbf{Consistency check.} In all three cases,
$\Vz+\Vk=\Vs$. These are analytical predictions only;
component identification and experimental verification require measured data.
\end{document}
